% ============================================================================
%  SECTION 6 — FISTA : FAST ISTA
%  >>> A REMPLIR PAR LE MEMBRE 3 <<<
%
%  Depend des sections 4 (prox) et 5 (ISTA).
%
%  Contenu attendu (voir REPARTITION_DETAILLEE.md) :
%    - Le momentum de Nesterov : d'ou vient le point extrapole y_k.
%    - Iteration :
%        x_k     = \softthr( y_k - t \nabla f(y_k), t\lambda )
%        t_{k+1} = (1 + sqrt(1 + 4 t_k^2)) / 2
%        y_{k+1} = x_k + ((t_k - 1)/t_{k+1}) (x_k - x_{k-1})
%    - PREUVE (ou esquisse detaillee) DE LA CONVERGENCE EN O(1/k^2)
%      (suite d'energie de Beck & Teboulle, 2009).
%    - Discussion des OSCILLATIONS / non-monotonie ; variantes "restart".
%
%  Renvoyer a la figure de convergence (section 9) qui montre O(1/k) vs O(1/k^2).
% ============================================================================
\section{FISTA : acceleration par momentum}
\label{sec:fista}

% --- [MEMBRE 3] Ecrire la section ici. Squelette suggere ci-dessous. --------

\subsection{Du gradient proximal au momentum de Nesterov}
% TODO (Membre 3)

\subsection{L'algorithme}
% TODO (Membre 3) : pseudo-code, une ligne de plus qu'ISTA.

\subsection{Convergence en $O(1/k^2)$}
% TODO (Membre 3) : enonce + preuve/esquisse.

\subsection{Non-monotonie et variantes a restart}
% TODO (Membre 3) : commenter les oscillations observees experimentalement.
